\section{Introduction}
\label{sec:intro}

Trading is increasingly delegated to algorithms that learn from their own
profits. A growing literature shows that such algorithms can arrive at
\emph{supracompetitive} outcomes without communicating: reinforcement-learning
(RL) pricing agents settle on prices above the competitive level
\citep{calvano2020,klein2021}, algorithmic market makers quote wider spreads
\citep{colliard2026,cont2024}, and RL informed speculators trade less
aggressively than competition would imply, earning collusive profits and
degrading price informativeness \citep{dou2025}. For regulators the finding is
alarming, because tacit collusion by algorithms falls outside the reach of
rules written around agreements among humans \citep{harrington2018,ezrachi2016}.

A supracompetitive outcome, however, is not the same thing as collusion. In the
game-theoretic sense, collusion is an outcome sustained by the threat of
punishment: a trader who deviates to grab a larger share of the rent is
punished by its rivals, so deviating does not pay. The same low-intensity
outcome can instead come from a learning bias: an algorithm that systematically
undervalues aggressive actions will trade conservatively even when no rival
could ever punish it. \citet{dou2025} name both channels in the trading context
(``price-trigger'' collusion and ``over-pruning'' bias), and a parallel debate
in industrial organization asks whether Q-learning collusion is genuine or an
artifact of exploration and learning design
\citep{calvano2023,abada2023,banchio2023,asker2022,denboer2026}. The
distinction is not academic. If algorithms punish, transparency matters,
because monitoring is what makes punishment possible, and the natural remedies
target information: reporting delays, reduced order-flow disclosure. If the
outcome is a learning bias, transparency is irrelevant and the remedies target
algorithm design.

This paper asks how to tell the two apart, and what the answer is for RL
informed traders in the canonical \citet{kyle1985} market and in the
\citet{dou2025} extension with information-insensitive investors.

\paragraph{Contributions.}
\begin{enumerate}
\item \textbf{Diagnostics that can fail.} We combine a paired
  \emph{rival-deviation} test (one trader is forced to play its one-period
  best response; we measure the rivals' reaction under common random numbers)
  and the \citet{dou2025} \emph{noise-shock} test with three placebos and
  falsification tests: a \emph{myopic placebo} (the same learners with
  $\gamma=0$, which have no reason to punish or fear punishment), an
  \emph{uninformative memory} (a random state of the same size as an
  informative one), and \emph{passive traders} that move the collusive outcome
  above the Nash one, so that collusion and bias predict opposite directions.
\item \textbf{A detectability result.} In the standard Kyle market a
  one-period deviation from the collusive intensity moves order flow by only
  $(I-1)/(2I)$ noise standard deviations per unit of $v/\sigma_v$, whatever the
  volume of noise trading (Proposition~\ref{prop:snr}). With a large mass of
  information-insensitive investors the ratio becomes
  $\xi\sigma_v(I-1)/(4I\sigma_u)$, about 625 in the \citet{dou2025}
  calibration. This explains from first principles why punishment-based
  collusion has been reported only in that regime.
\item \textbf{The bias, isolated.} A single informed trader with no rival and
  a fixed price rule under-trades by up to 57\% under standard Q-learning, by
  an amount proportional to $\sqrt\alpha$ ($R^2=\MechSqrtRsq$), while
  counterfactual updating learns the optimum.
\item \textbf{Low trading without collusion.} In the standard Kyle market,
  forward-looking Q-learners close on average \ResidDeltaInt{} of the gap
  between competition and the cartel, yet rivals never punish deviations,
  even with perfect monitoring, and deviating pays. The myopic placebo trades
  even less ($\Delta=\MyopicResidDeltaInt$, beyond the cartel); an uninformative
  memory reproduces the outcome ($\RandomMemDeltaInt$); and when passive
  traders put the cartel \emph{above} Nash, learners still trade half their
  Nash intensity and earn \PassiveProfitPct\% less than in equilibrium.
\item \textbf{What fixes it.} Blurring the order flow traders observe changes
  nothing ($\Delta=\OpaqueDeltaInt$), while counterfactual updating, a change
  to the learning rule alone, restores competition ($\Delta=\CfResidDeltaInt$).
  In the \citet{dou2025} regime, where deviations are glaring, our
  separate-table learners match the independent replication's collusion
  level ($\DouStratDeltaInt$) and still do not punish. With a shared Q-table, a \emph{myopic} learner
  reproduces both the high collusion index ($\SharedDouMyopicDeltaInt$) and the
  threshold-shaped response to noise shocks that has been read as a price
  trigger, so that signature does not identify collusion without a placebo.
\item \textbf{An open testbed.} A batched, CPU-only simulator with exact
  benchmarks, tabular Q-learning, DQN and PPO agents, all diagnostics, and a
  test suite that validates the economics and the diagnostics against
  scripted traders whose behaviour is known.
\end{enumerate}

\paragraph{Why this matters for policy.} The two channels call for different
interventions. Our results suggest that in markets like the standard Kyle
setting, the supracompetitive outcomes produced by off-the-shelf Q-learning are
properties of the learning procedure, and would respond to algorithm-design
standards (step sizes, exploration) rather than to limits on transparency.
Where deviations are highly detectable, as with a large mass of
information-insensitive investors, the diagnostics must be applied with a
placebo, because a reaction to a price shock is not by itself evidence of
punishment.
